% Body-only proof section for the final 244-state construction.
\section{The transfer-plus-one rotation}
\label{sec:qt-rotation}

The useful reduction is a loop rotation in the proof enumerator.  It is
valid for every register vector and needs no reachability invariant.
The source already shares its transfer loop with its following increment;
moving that increment to the beginning of each iteration allows the
three-state local decoder to become a two-state decoder.  The complete
runtime certificates are needed to confirm this saving because the
machine also reuses source decoders as phase routers.

\subsection{The exact source change}

Let $N=R_0$ be the next-proof register and $X=R_6$ its scratch register.
The old closed fragment has the following behavior:
\begin{quote}
While $X>0$, decrement $X$ and increment $N$.  After the loop,
increment $N$ once more.
\end{quote}
The replacement performs the increment before the test:
\begin{quote}
Increment $N$.  If $X=0$, exit.  Otherwise decrement $X$ and repeat.
\end{quote}
The zero test here is a natural-register test: it leaves zero unchanged.
This avoids any convention about signed arithmetic or a programming
language's postfix decrement.

In the literal source tree the fragment is rooted at
$p=\texttt{011}$.  Its old and new leaf paths are
\[
\begin{array}{c|c|c@{\qquad}c|c|c}
\multicolumn{3}{c}{\text{Old fragment}}&
\multicolumn{3}{c}{\text{New fragment}}\\
\text{site}&\text{path}&\text{instruction}&
\text{site}&\text{path}&\text{instruction}\\\hline
6&p00&\operatorname{dec}X&6&p0&\operatorname{inc}N\\
7&p010&\operatorname{inc}N&7&p10&\operatorname{dec}X\\
8&p011&\operatorname{continue}(2)&8&p11&\operatorname{continue}(2)\\
9&p1&\operatorname{inc}N&&&
\end{array}
\]
Thus four source leaves become three.  The complete program has 6,459
source leaves and 8,390 guarded instruction cases.  These numbers count
source instructions and proof obligations, not Turing-machine states.

Writing $\mathsf{Node}$ for an ordered binary source node, the change is
\[
\begin{split}
&\mathsf{Node}\bigl(
 \mathsf{Node}(\operatorname{dec}X,
   \mathsf{Node}(\operatorname{inc}N,\operatorname{continue}(2))),
 \operatorname{inc}N\bigr)\\
&\hspace{2em}\longmapsto
 \mathsf{Node}\bigl(\operatorname{inc}N,
   \mathsf{Node}(\operatorname{dec}X,\operatorname{continue}(2))\bigr).
\end{split}
\]
This is a proved loop transformation, rather than a claim that arbitrary
reassociation of break/continue trees preserves behavior.

\subsection{A universal arithmetic proof}

Suppose the entry values are $N=a$ and $X=b$, where $a,b\in\mathbb N$.
Both fragments terminate with
\[
 N=a+b+1,\qquad X=0,
\]
and leave the other seven registers unchanged.  In fact, both execute
exactly $3b+2$ source instructions.

For $b=0$, the old fragment takes its zero branch and then increments
$N$.  The new fragment increments $N$ and then takes its zero branch.
They each execute two instructions and return the same register vector.

For $b>0$, the old fragment executes a positive decrement of $X$, an
increment of $N$, and a continue.  The new fragment executes the same
increment and decrement in the opposite order, followed by a continue.
The two arithmetic operations affect distinct registers and commute.
After three instructions, both fragments return to their respective
entry sites with $(N,X)=(a+1,b-1)$.  Induction on $b$ proves both the
postcondition and the exact instruction count.  It also rules out an
infinite execution confined to either fragment.

The checker verifies these two cases symbolically over unbounded regions.
In the first case, $X=0$; in the second, $X\geq1$.  Each register is an
independent entry variable plus an integer offset.  The checker propagates
the bounds needed to force each branch and compares the final offsets.
The zero case gives $\Delta N=1$, and the positive case gives
$\Delta N=1$, $\Delta X=-1$.  It does not substitute a finite selection of
register values for this argument.

An equivalent relation can be placed at the decrement instructions:
at related old site 6 and new site 7, the new $N$ exceeds the old $N$ by
one and all other registers agree.  A positive iteration preserves this
relation.  The old fragment's final increment restores equality on exit.
The equal-register relation at the two fragment entries is sufficient for
the full context proof below.

\subsection{The complete source context}

Outside the four old fragment sites, old source site $i<6$ corresponds to
new site $i$, and old site $i\geq10$ corresponds to new site $i-1$.
The fragment entry is mapped from old site 6 to new site 6.  The checker
reconstructs both complete binary trees and verifies all 8,385 successor
edges outside the fragment, including their exact break and continue
paths and implicit zero padding.  The only external edge entering the
fragment is the zero branch of old site 2.  No external edge enters one
of its interior instructions.

Consequently runs may be compared at every outside instruction and at
every fragment entry.  Outside the fragment, one instruction on each
side has the same register effect and reaches corresponding sites.
Inside it, the preceding induction supplies finite matching executions
with identical exit registers.  The halt instruction is outside the
fragment: old site 6324 becomes new site 6323.  Therefore, for every
$r\in\mathbb N^9$,
\[
 \operatorname{Halts}(P_{245},r)
 \quad\Longleftrightarrow\quad
 \operatorname{Halts}(P_{\mathrm{rot}},r).
\]
The same argument preserves nonhalting.  There is no hidden divergence
inside a fragment, since $X$ strictly decreases on every positive
iteration and the zero case exits.

All ingredients are arithmetic on finite instruction tables and natural
registers, with induction on a natural number.  Thus this implication
and its converse are ordinary theorems formalizable in ZFC.  No
consistency assumption is used in proving the source equivalence.

\subsection{Why the decoder saves one working state}

In the old labels, the relevant decoder roles are 12, 17, and 26.
Role 12 selects the transfer loop or its final increment.  Role 17
selects the decrement or its two-instruction loop body, and role 26
selects the body increment or the continue operation.
The replacement uses role 12 to select the increment or the test/continue
pair, and role 17 for that pair.  Role 26 is deleted.  In the old table,
17 has only predecessor $(12,0)$ and 26 has only predecessor $(17,1)$;
the transfer loop has no additional caller requiring the old role 26.

The three changed retained transition rows are
\[
\begin{array}{c|c|c}
\text{read row}&\text{old transition}&\text{new transition}\\\hline
(12,0)&(0,R,17)&(0,R,18)\\
(12,1)&(1,R,18)&(1,R,17)\\
(17,1)&(1,R,26)&(0,L,20).
\end{array}
\]
Both rows of state 26 are removed, and every other row is unchanged.
The new first increment has path \texttt{0110}, ending in the same
two bits \texttt{10} as the old loop increment \texttt{011010}.
This observation suggests that the reused increment decoder is valid in
its new context.  The complete symbolic runtime certificates, rather
than that suffix observation alone, discharge the decoder and phase-router
obligations for every source site and register valuation.

\paragraph{Independent replay.}
The self-contained script \path{check_qt_rotation.py} imports no
compiler or inherited proof code.  It checks source-tree well-formedness,
every surrounding control edge, both unbounded arithmetic cases, and
the halt-site correspondence.  It also checks the alternative left
association
$\mathsf{Node}(\mathsf{Node}(\operatorname{inc}N,
\operatorname{dec}X),\operatorname{continue}(1))$ at the source level;
the delivered literal table uses the right association displayed above.
The JSON certificate records the exact old and new source hashes and
the symbolic branch traces.  This source proof is composed with the
separate blank-tape bootstrap and full runtime proof for the 244-state
table.
